\section{A Restricted Intent Compiler}
\label{sec:compiler}

An automated compiler cannot evaluate~\eqref{eq:intent}. Membership in
\(\mathrm{Sat}(\Phi_{\text{func}})\) is undecidable for interesting
contracts, by Rice's theorem~\cite{rice1953classes}, and even a decidable
proof still has a time budget. This section replaces the source
step~\eqref{eq:intent_source} with a generator, a sound verifier \(V\) with
budget \(t\), and a fallback. The as-if step~\eqref{eq:intent_lower} remains
\(\mathcal{L}_{\text{Target}}\). The contract bounds what the compiler may
do, and what \(V\) proves within \(t\) bounds what it can use.

\subsection{Generator}
\label{sec:generator}

Let \(R\) be a catalog of source-to-source rewrites, each with a syntactic
applicability condition and optional side constraints on parameters, for
example that capacity is a power of two. Let \(D\) be a finite domain for
parameter holes, as machine integers are. \(R\) and \(D\) are inputs of
\(\mathrm{Gen}\). The \emph{generated set} is
\[
\mathrm{Gen}(P_1, R, D)
= \{ S[h] \mid S \text{ is obtained from \(P_1\) by a finite sequence of
rewrites from \(R\)},\ h \in D \}.
\]
A program is in \(\mathrm{Gen}\) only if a finite sequence of rewrites from
\(R\) produces its structure and some \(h \in D\) fills its holes, so a
cursor-based PID allocator, a linked queue, or a ring that pads storage
while keeping caller-facing capacity appears only if \(R\) contains that
rewrite. Rewrites may change observable behavior, and their outputs are
candidates until verified.

\subsection{Verifier}
\label{sec:verifier}

A verifier \(V\) takes a program, a contract, and a time budget \(t\) and
answers \(\mathsf{proved}\), a counterexample, or \(\mathsf{unknown}\). It
is \emph{sound} if \(\mathsf{proved}\) implies
\(P \models \Phi_{\text{func}}\) under an exact model of machine arithmetic,
including overflow.

Many useful verifiers are bounded. A bound \(K\) (operations, unwindings)
induces a judgment \(P \models_K \Phi_{\text{func}}\): the contract holds on
every execution the harness covers up to \(K\), and a sound bounded verifier
may conclude \(\mathsf{proved}\) only for that judgment. Timeouts, solver
failures, exhausted unwindings, and a property failure on a truncated path
are \(\mathsf{unknown}\), never \(\mathsf{proved}\).

There are two ways to certify a candidate \(P_2\). A relational check proves
that \(P_2\) behaves like \(P_1\), taking observational equivalence to
\(P_1\) as the bound. The
developer writes no \(\Phi_{\text{func}}\) and the query contains both
programs. Every candidate outside \([P_1]_\approx\) is then rejected.
That query inherits
\(P_1\)'s complexity and may be undecidable, but it is cheap when the
programs are structurally close. Those checks remain the right tool for a
small patch when no contract is available~\cite{lopes2021alive2}. A
decoupled check proves \(P_2 \models \Phi_{\text{func}}\) (or
\(\models_K\)) directly, taking \(\Phi_{\text{func}}\) as the
specification. The developer must write that contract, and the query
contains only the candidate. This check does not inherit \(P_1\)'s complexity, but the
candidate itself has to lie in a fragment \(V\) can decide, or decide up to
\(K\). \(P_1\) need not satisfy the contract: it is a starting point for the
generator and a fallback. Search stays near that baseline so a candidate is
still a program the developer can read against the sources they already
have.

This compiler uses the decoupled check, so the \emph{admitted set} at budget
\(t\) is
\[
\mathcal{A}(P_1, \Phi, t)
= \{ P_2 \in \mathrm{Gen}(P_1, R, D) \mid
V(P_2, \Phi_{\text{func}}, t) = \mathsf{proved}
\ \text{and}\ P_2 \text{ passes the check for }\mathcal{B} \}.
\]
The two conjuncts are different checks: \(V\) admits on \(\Phi_{\text{func}}\),
and a provided checker admits on \(\mathcal{B}\), not a second run of \(V\).
When \(V\) is unbounded and sound, \(\mathcal{A}\) is a subset of
\(\mathrm{Sat}(\Phi_{\text{func}})\). When \(V\) is bounded, it is a subset
of the \(K\)-satisfying programs. Only programs in \(\mathrm{Gen}\) are
submitted to \(V\), and \(\mathsf{unknown}\) shrinks \(\mathcal{A}\) without
unsounding an emit. A
bounded model checker or another sound oracle can play \(V\), and they
differ in which contracts they express and how much of \(\mathrm{Gen}\) they
reach in \(t\). A synthesizer built on \(V\) can produce only what \(V\) can
prove~\cite{Srivastava2010VS3}.

\subsection{Backend}
\label{sec:backend}

The as-if step~\eqref{eq:intent_lower} is unchanged: the restricted compiler
hands a source program to \(\mathcal{L}_{\text{Target}}\).

\paragraph{Compiler Correctness Assumption.}
For every program \(P\) and every \(\sigma\) on which \(P\) has no undefined
behavior, \(\llbracket \mathcal{L}_{\text{Target}}(P) \rrbracket(\sigma)\)
and \(\llbracket P \rrbracket(\sigma)\) have the same observable part. That
is the as-if rule as compilers implement it, trusted rather than verified.

Proposition~\ref{prop:refine} already says that an as-if program with the
same \(\theta\) remains legal. \(\mathcal{L}_{\text{Target}}\) is that as-if
on inputs where the source has no undefined behavior, and it keeps
\(\theta\), so the unbounded case of the next lemma is that proposition
instantiated at \(P' = \mathcal{L}_{\text{Target}}(P^\star)\). The bounded
judgment is the same fact on the traces \(\models_K\) covers.

\begin{lemma}[Lowering Preserves a Proved Contract]
\label{lem:lowering}
If \(P^* \models \Phi_{\text{func}}\) (respectively \(\models_K\)), then
\(\mathcal{L}_{\text{Target}}(P^*) \models \Phi_{\text{func}}\)
(respectively \(\models_K\)).
\end{lemma}

\noindent\emph{Proof sketch.} Implementation parameters remain those of the
source. \(P^*\) has no undefined behavior on inputs in \(\mathit{Pre}\) (or
on the traces \(\models_K\) covers), so the compiler correctness assumption
applies, and Proposition~\ref{prop:refine} gives the unbounded case.
(Appendix~\ref{app:proofs}.)

The lemma needs a contract over observable results and a candidate that is
UB-free on the covered inputs, so the trusted base is the production
compiler, the verifier, and the encoding of the language into the verifier's
logic, which is larger than that of a verified compiler or of Fiat's
end-to-end proofs~\cite{leroy2009backend,delaware2015fiat}.

\subsection{Selection and Fallback}
\label{sec:selection}

Winner selection lowers the baseline and every admitted candidate and keeps
the binary with the lowest \emph{measured} cost:
\begin{equation}
\label{eq:winner}
\mathrm{Bin}^* = \arg\min_{\mathrm{Bin} \in \{\mathcal{L}_{\text{Target}}(P)
\mid P \in \mathcal{A} \cup \{P_1\}\}}
\mathcal{O}(\mathrm{Bin}, \text{Target}),
\end{equation}
preferring \(P_1\) on ties. Each admitted source, and \(P_1\), is compiled
with \(\mathcal{L}_{\text{Target}}\) and ranked by measured cost on the
profiled workload and hardware, and the emitted artifact is the winning
source, not an object file.

If \(\mathcal{A}\) is empty, the compiler emits \(P_1\) unchanged. A change
is admitted only after a proof of \(\Phi_{\text{func}}\) (or \(\models_K\)).
Callers whose tests depended on incidental behavior the contract left open
can still observe a difference.

\subsection{What the Restricted Compiler Guarantees}
\label{sec:restricted_thm}

\begin{theorem}[Soundness of a Restricted Intent Compiler]
\label{thm:restricted}
Let \(V\) be sound for \(\models\) or for \(\models_K\), and let the compiler
emit source \(P^*\) and binary
\(\mathrm{Bin}^* = \mathcal{L}_{\text{Target}}(P^*)\).
\begin{enumerate}
  \item If \(P^* \neq P_1\), then \(P^*\) is in \(\mathrm{Gen}(P_1, R, D)\),
  \(V\) proved it, and it passed \(\mathcal{B}\). Hence
  \(P^* \models \Phi_{\text{func}}\) or \(P^* \models_K \Phi_{\text{func}}\),
  matching \(V\), and the same judgment holds of \(\mathrm{Bin}^*\) by
  Lemma~\ref{lem:lowering}.
  \item The measured cost of \(\mathrm{Bin}^*\) is at most that of
  \(\mathcal{L}_{\text{Target}}(P_1)\), when measurement runs.
  \item If \(P_1 \models \Phi_{\text{func}}\) (or \(\models_K\)), then
  \(\mathrm{Bin}^*\) satisfies the same judgment in every case.
\end{enumerate}
\end{theorem}

\noindent\emph{Proof sketch.} Winner selection returns \(P_1\) or a member
of \(\mathcal{A}\). Membership in \(\mathcal{A}\) is part 1. Part 2 is
minimization over a set that contains the compiled baseline. Part 3 uses
part 1 or Lemma~\ref{lem:lowering} on \(P_1\).
(Appendix~\ref{app:proofs}.)

If \(P_1\) already satisfies the contract, every emit does. If it does not,
the compiler replaces it only with a proved program or leaves it as written,
so the emit is never less safe than compiling \(P_1\) with
\(\mathcal{L}_{\text{Target}}\) alone. If some generated hole assignment of
a given sketch satisfies the contract and \(V\) decides every query, a
baseline-directed search that enumerates \(D\) in order of distance to
\(\theta_{P_1}\) returns a closest satisfying assignment
(Theorem~\ref{thm:cegis}). If no generated program satisfies the contract,
the compiler keeps \(P_1\), and a different structure requires a rewrite in
\(R\).

The compiler ranks by \(\mathcal{O}\) among programs \(V\) has proved.
How \(\mathrm{Gen}\) is explored, including any score that orders rewrites,
does not change that ranking. Verification cost and execution cost need not
agree~\cite{wagner2013overify,rinard2026quantitative}, so the fastest program
may lie outside \(\mathcal{A}\). They agree on the ring: integer division is
expensive on the processor and expensive to
bit-blast~\cite{niemetz2024bitblasting}, and the rewrite that removes the
remainder also simplifies the proof.
